% ============================================================================
\chapter{同余子群}
\label{ch:congruence-subgroups}
% ============================================================================

% ----------------------------------------------------------------------------
\section{$\GammaO(N)$、$\GammaI(N)$ 和 $\Gamma(N)$ 的定义}
% ----------------------------------------------------------------------------

\begin{definition}[同余子群]
设 $N$ 为正整数。定义以下 $\SL_2(\Z)$ 的子群：
\begin{align}
  \Gamma(N) &= \left\{
    \begin{pmatrix} a & b \\ c & d \end{pmatrix} \in \SL_2(\Z)
    \;\middle|\;
    \begin{pmatrix} a & b \\ c & d \end{pmatrix}
    \equiv \begin{pmatrix} 1 & 0 \\ 0 & 1 \end{pmatrix} \pmod{N}
  \right\}, \\[6pt]
  \GammaI(N) &= \left\{
    \begin{pmatrix} a & b \\ c & d \end{pmatrix} \in \SL_2(\Z)
    \;\middle|\;
    \begin{pmatrix} a & b \\ c & d \end{pmatrix}
    \equiv \begin{pmatrix} 1 & * \\ 0 & 1 \end{pmatrix} \pmod{N}
  \right\}, \\[6pt]
  \GammaO(N) &= \left\{
    \begin{pmatrix} a & b \\ c & d \end{pmatrix} \in \SL_2(\Z)
    \;\middle|\;
    c \equiv 0 \pmod{N}
  \right\}.
\end{align}
包含关系为 $\Gamma(N) \trianglelefteq \GammaI(N) \trianglelefteq \GammaO(N)
\leq \SL_2(\Z)$。
\end{definition}

% ----------------------------------------------------------------------------
\section{指标与陪集分解}
% ----------------------------------------------------------------------------

% TODO: 各子群在 SL₂(ℤ) 中的指标公式

% ----------------------------------------------------------------------------
\section{基本域与尖点}
% ----------------------------------------------------------------------------

% TODO: 同余子群的基本域、尖点的分类

% ----------------------------------------------------------------------------
\section{习题}
% ----------------------------------------------------------------------------
